\Qhead{3.8}{Случайные поля. Когерентность и дифракция}{Случайные поля. Пространственная и временная когерентность. Дифракция случайных волн. Теорема Ван Циттерта-Цернике. Дифракция регулярной волны на случайном фазовом экране. Тепловое электромагнитное поле. Теорема взаимности.}
\Subfile{3.8.1}{sections/03_statistics/q08/s01}
\Subfile{3.8.2}{sections/03_statistics/q08/s02}

\recap{Когерентность — корреляция полей, измеряемая интерференцией. Дифракция переносит взаимную когерентность по двум координатам. Некогерентный источник создаёт частично когерентное поле; фазовый экран разрушает корреляции через разность фаз.}
\board{Видность полос, двойной интеграл переноса, преобразование Фурье яркости, $\exp(-D_\phi/2)$ и условия взаимности.}
\oral{Почему равенство нулю среднего поля не означает отсутствия света?}
\bait{Можно оставить тепловое группирование: гауссовское спаривание даёт удвоенную интенсивностную корреляцию при нулевой задержке.}
\sources{Barbastathis~\cite{barbastathis}, материалы Optical Engineering о статистической оптике; Fang~\cite{fangcoherence}, обзор пространственной когерентности. Формулы экрана, Зигерта и взаимности выведены здесь при явно указанных предположениях.}
